\section{Conos proyectivos, ejemplo de Zariski}
\label{XI.3}
\marginpar{290}

Sea, como antes y para simplificar, $k$ un cuerpo algebraicamente cerrado,
y sea $V$ un $k$-esquema proyectivo conexo, subesquema cerrado
de~$\PP^r_k$, que, si se desea, puede suponerse no singular. Sean
$Y=\hat C$ el cono proyectivo sobre~$V$, $y_0$ su vértice,
$X=\hat C_V$ la clausura proyectiva usual del fibrado vectorial
$C_V=\VV(\cal{O}_V(1))$ asociado a~$\cal{O}_V(1)$ y, por último,
$$
f\colon X\to Y
$$
el morfismo canónico que contrae a un punto la sección nula $X_0$
de~$C_V$ en~$X$ (EGA~II~8.6.4). Como $X$ es un fibrado localmente trivial
sobre~$V$, con fibras $\PP^1$ y, por tanto, simplemente conexas, el
morfismo $p\colon X\to V$ induce, en virtud
\ifthenelse{\boolean{orig}}
{de~(\Ref{XI.4}.)}
{de~(\Ref{X.1.4})}
un isomorfismo:
$$
\pi_1(X)\isomto\pi_1(V)\quoi.
$$
Como $p$ induce un isomorfismo $X_0\to V$, se concluye que
\emph{un recubrimiento étale de~$X$ está completamente descompuesto si y
solo si su restricción a~$X_0$ lo está}. Ahora bien, para todo recubrimiento
étale $Y'$ de~$Y$, la imagen inversa $X'=X\times_Y Y'$ es un recubrimiento
étale de~$X$ completamente descompuesto sobre la fibra~$X_0$ y, por tanto,
trivial. Como el homomorfismo $\pi_1(X)\to\pi_1(Y)$ es sobreyectivo
(\Ref{IX.3.4}), se concluye que
$$
\pi_1(Y)=(e)
$$
en otras palabras, \emph{todo cono proyectivo es simplemente conexo.}
(N.B.: en característica~$0$, el mismo resultado sigue siendo válido si se
toma como~$Y$ el cono afín.)

Supongamos ahora que $V$ es regular, \ie lisa sobre~$k$; entonces $X$ es
regular y, para una inmersión proyectiva conveniente de~$V$, se obtiene un
cono proyectivo~$Y$ \emph{normal}. Si $V$ no es simplemente conexa, y por
tanto tampoco lo es~$X$, sea $X'$ un recubrimiento étale conexo no trivial
de~$X$. Como las fibras de~$X$ sobre los puntos $y\in Y$ distintos de~$y_0$
se reducen a un punto, se ve que la restricción de~$X'$ a estas fibras
(en particular, a la fibra
\marginpar{291}
genérica) es trivial; sin embargo,
\ifthenelse{\boolean{orig}}
{$X$}
{$X'$}
no proviene por imagen inversa de un recubrimiento étale de~$Y$, pues $Y$
es simplemente conexa y $X'$ estaría completamente descompuesto. Esto
muestra que (\Ref{X.1.3} y~\Ref{X.1.4}) se vuelven falsos si se sustituye
la hipótesis de que $f$ es separable por la más débil de que sus fibras son
esquemas algebraicos separables (o incluso lisos) sobre los~$\kres(s)$.
Obsérvese asimismo que los grupos fundamentales de las fibras geométricas
$\overline X_y$ para $y\neq y_0$ se reducen evidentemente a~$(e)$, puesto
que esas fibras se reducen a un punto, mientras que $\pi_1(X_0)\neq e$;
por tanto, el teorema de semicontinuidad (\Ref{X.2.4}) también falla
para~$f$.

Indiquemos, por último, el ejemplo señalado por Zariski que hace fallar
estos mismos teoremas cuando se sustituye en ellos la hipótesis de que $f$
es separable por la de que $f$ es plano. Sea $f\colon X\to Y$ un morfismo
de una superficie proyectiva no singular a la recta racional $Y=\PP^1$,
tal que
\ifthenelse{\boolean{orig}}
{$K=k(x)$}
{$K=k(X)$}
sea una extensión «regular», \ie primaria y separable, de~$k(f)$ (\ie la
fibra genérica geométrica es conexa y separable), y tal que el divisor
$(f)=X_0-X_\infty$ sea
\ifthenelse{\boolean{orig}}
{un múltiplo $n$.º de un divisor}
{$n$ veces un divisor}
(donde $n$ es un entero primo con la característica). Es posible construir
tales ejemplos en toda característica. Sea $X'$ la normalización de~$X$ en
$K(f^{1/n})$, donde $K=k(X)$ es el cuerpo de funciones de~$X$. De la
hipótesis sobre~$(f)$ resulta que $X'$ es étale sobre~$X$. Sea $Y'$ el
normalización de~$Y$ en $k(t)(t^{1/n})$; está ramificada sobre~$Y$ exactamente
en los puntos $t=0$ y $t=\infty$,
\ifthenelse{\boolean{orig}}
{y la restricción $X'|f^{-1}(U)$ es isomorfa a la imagen inversa de~$Y'|U$.}
{y, poniendo $U=Y-\{0,\infty\}$, la restricción $X'|f^{-1}(U)$ es isomorfa a la imagen inversa de~$Y'|U$.}
En particular, la restricción de~$X'$ a la fibra genérica \emph{geométrica}
de~$X$ sobre~$Y$ se descompone completamente. Sin embargo, $X'$ no es
isomorfo a la imagen inversa de un recubrimiento étale de~$Y$, pues se ve
enseguida que este último sería necesariamente~$Y'$, lo cual es absurdo
porque $Y'$ está ramificado sobre~$Y$\kern2pt\footnote{Puede observarse,
desde el punto de vista de la «topología étale» (SGA~4~VII), que en este
ejemplo $\R^1(f_{\et})_*(\ZZ/n\ZZ)$ es «no separado» sobre
\ifthenelse{\boolean{orig}}{$S$}{$Y$}.
}.

He aquí (siguiendo a Serre) una manera sencilla de realizar las condiciones
de este ejemplo, inspirada en~\cite[\No 20]{XI.5}: se toma como~$n$ un
número primo $\geq5$, distinto de la característica, y se hace actuar
$G=\ZZ/n\ZZ$ sobre $\underline{k}^4$%
\ifthenelse{\boolean{orig}}
{}
{{\renewcommand{\thefootnote}{*}\addtocounter{footnote}{-1}\kern1pt\footnote{\lcrochetbf Añadido en 2003: $\underline{k}^4$ designa el espacio afín de dimensión~$4$ sobre~$k$.\rcrochetbf}}}
multiplicando las coordenadas por cuatro caracteres distintos de~$G$ (lo
cual es posible porque $n\geq5$). Entonces $G$ actúa sobre el espacio
proyectivo $\PP^3_k$, y los únicos puntos fijos por~$G$ son los cuatro
puntos correspondientes a los ejes coordenados. La superficie $X'$ de
ecuación
\ifthenelse{\boolean{orig}}
{$x^4+y^4+z^4+t^4=0$}
{$x^n+y^n+z^n+t^n=0$}
es lisa sobre
\marginpar{292}
$k$ (criterio jacobiano) y no contiene ninguno de los puntos fijos; por
tanto, como $G$ es de orden primo,
\ifthenelse{\boolean{orig}}
{actúa sobre~$X$ «sin puntos fijos», \ie $X$ es un recubrimiento principal de~$X=X'/G$ de grupo~$G$.}
{actúa sobre~$X'$ «sin puntos fijos», \ie $X'$ es un recubrimiento principal de~$X=X'/G$ de grupo~$G$.}
Sea $g=x/y$ en $k(X')=K'$; es un generador de Kummer de~$K'$ sobre
$K=k(X)$ si los caracteres elegidos eran~$\chi^i$, $i=0,1,2,3$, con
$\chi$ un carácter primitivo. Sea $f$ su potencia $n$-ésima, que es un
elemento de~$K$. Se ve enseguida que $K'$ es una extensión regular
de~$k(g)$; esto resulta de que la curva plana de ecuación homogénea en
$U,T,Z$,
$T^n+Z^n+(1+g^n)U^n=0$, es lisa sobre~$k(g)$ (criterio jacobiano), y de que
se sabe que toda curva plana es conexa. Por otra parte,
$k(f)=K\cap k(g)$, pues el segundo miembro es una extensión de~$k(f)$
contenida en la extensión $k(g)$ de grado primo y
\ifthenelse{\boolean{orig}}{distinto}{distinto}
de~$k(g)$ (pues $g\not\in K$). Esto implica que $K$ es una extensión regular
de~$k(f)$. Finalmente, el divisor de~$f$ sobre~$X$ es $n$ veces un divisor,
pues su imagen inversa sobre~$X'$ es el divisor de~$g^n$, luego es $n$
veces un divisor, y puede hacerse el descenso porque $X'$ es étale
sobre~$X$. Se habría terminado si la aplicación racional
$f\colon X\to\PP^1$ fuese un morfismo, es decir, si los divisores de los
ceros y de los polos de~$f$ no se intersecaran. De hecho, se comprueba
fácilmente (mirando de nuevo sobre~$X'$) que los dos divisores en cuestión
son respectivamente $n$ veces sendas curvas lisas sobre~$k$ que se cortan transversalmente en
un punto~$a$. Sustituyendo ahora $X$ por el esquema obtenido mediante una
explosión en~$a$, siguen verificándose las condiciones anteriores
($\divisor(f)$ divisible por~$n$, y $k(X_1)=k(X)$ extensión regular de
\ifthenelse{\boolean{orig}}{$k(X)$}{$k(f)$})
pero, además, $f$ es un \emph{morfismo} $X_1\to\PP^1$; por tanto, se está
en las condiciones buscadas.
